% ============================================================
\chapter{AR, MA, ARMA Models}
\label{ch:arma}
% ============================================================

\section{Motivating example}
The daily log-returns of the CAC~40 exhibit weak but significant correlation at the first few lags. An AR(1) or ARMA(1,1) model captures this linear dependence.

\section{Autoregressive processes AR$(p)$}\index{AR}

\begin{definition}[AR$(p)$]
\[
X_t = \phi_1 X_{t-1} + \phi_2 X_{t-2} + \cdots + \phi_p X_{t-p} + \eps_t,
\]
where $\eps_t\sim\mathrm{WN}(0,\sigma^2)$. In operator notation: $\phi(B)X_t = \eps_t$ with $\phi(z)=1-\phi_1 z-\cdots-\phi_p z^p$.
\end{definition}

\begin{theorem}[Stationarity of an AR$(p)$]\label{thm:ar-stat}
The AR$(p)$ admits a unique stationary solution if and only if all roots of $\phi(z)=0$ lie \textbf{outside the unit disc}: $|z_i|>1$ for all $i$.
\end{theorem}

\begin{proof}[Sketch for AR(1)]
$X_t = \phi X_{t-1}+\eps_t$. By recursion: $X_t = \sum_{j=0}^{\infty}\phi^j\eps_{t-j}$. This series converges in $L^2$ if and only if $\sum\phi^{2j}<\infty$, i.e.\ $|\phi|<1$. The root of $\phi(z)=1-\phi z$ is $z=1/\phi$, which lies outside the disc if and only if $|\phi|<1$.
\end{proof}

\begin{theorem}[ACF of an AR(1)]
If $X_t = \phi X_{t-1}+\eps_t$ with $|\phi|<1$:
\[
\gamma(h) = \frac{\sigma^2\phi^{|h|}}{1-\phi^2}, \quad \rho(h)=\phi^{|h|}.
\]
The ACF decays exponentially.
\end{theorem}

\begin{proof}
$\gamma(h)=\E[X_t X_{t+h}]=\phi\gamma(h-1)$ for $h\geq 1$ (by multiplying the AR equation by $X_{t+h}$ and taking expectations). With $\gamma(0)=\sigma^2/(1-\phi^2)$.
\end{proof}

\begin{theorem}[PACF of an AR$(p)$]
For an AR$(p)$: $\alpha(h)=0$ for all $h>p$. The PACF truncates at order $p$.
\end{theorem}

\section{Yule--Walker equations}\index{Yule--Walker}

\begin{definition}
For a stationary AR$(p)$, the relations:
\[
\gamma(h) = \phi_1\gamma(h-1)+\cdots+\phi_p\gamma(h-p) \quad (h\geq 1)
\]
yield the \textbf{Yule--Walker} system:
\[
\begin{pmatrix}\gamma(0)&\gamma(1)&\cdots&\gamma(p-1)\\\gamma(1)&\gamma(0)&\cdots&\gamma(p-2)\\\vdots&&\ddots&\vdots\\\gamma(p-1)&\gamma(p-2)&\cdots&\gamma(0)\end{pmatrix}
\begin{pmatrix}\phi_1\\\phi_2\\\vdots\\\phi_p\end{pmatrix}
=\begin{pmatrix}\gamma(1)\\\gamma(2)\\\vdots\\\gamma(p)\end{pmatrix}.
\]
\end{definition}

\section{Moving average processes MA$(q)$}\index{MA}

\begin{definition}[MA$(q)$]
\[
X_t = \eps_t + \theta_1\eps_{t-1} + \cdots + \theta_q\eps_{t-q} = \theta(B)\eps_t,
\]
where $\theta(z)=1+\theta_1 z+\cdots+\theta_q z^q$.
\end{definition}

\begin{theorem}[Stationarity of an MA$(q)$]
Any MA$(q)$ is \textbf{always stationary} (finite sum of white noise terms).
\end{theorem}

\begin{theorem}[ACF of an MA$(q)$]
\[
\gamma(h) = \begin{cases}\sigma^2\sum_{j=0}^{q-|h|}\theta_j\theta_{j+|h|} & \text{if }|h|\leq q,\\0 & \text{if }|h|>q,\end{cases}
\]
with $\theta_0=1$. The ACF \textbf{truncates} at lag $q$.
\end{theorem}

\begin{definition}[Invertibility]\index{invertibility}
An MA$(q)$ is \textbf{invertible} if all roots of $\theta(z)=0$ lie outside the unit disc. In that case, $X_t$ admits an AR$(\infty)$ representation: $\pi(B)X_t=\eps_t$ with $\pi(z)=\theta(z)^{-1}$.
\end{definition}

\section{ARMA$(p,q)$ processes}\index{ARMA}

\begin{definition}[ARMA$(p,q)$]
\[
\phi(B)X_t = \theta(B)\eps_t, \quad\text{i.e.}\quad X_t - \sum_{i=1}^p\phi_i X_{t-i} = \eps_t + \sum_{j=1}^q\theta_j\eps_{t-j}.
\]
\end{definition}

\begin{theorem}[Stationarity and invertibility conditions]
\begin{itemize}
\item Stationary $\iff$ roots of $\phi(z)=0$ outside the unit disc.
\item Invertible $\iff$ roots of $\theta(z)=0$ outside the unit disc.
\item One always requires that $\phi(z)$ and $\theta(z)$ share \textbf{no common root} (otherwise the model is over-parameterised).
\end{itemize}
\end{theorem}

\begin{theorem}[MA$(\infty)$ representation]
If the ARMA$(p,q)$ is stationary: $X_t = \psi(B)\eps_t$ with $\psi(z)=\theta(z)/\phi(z)=\sum_{j=0}^{\infty}\psi_j z^j$. The $\psi_j$ are the \textbf{impulse response weights}.
\end{theorem}

\section{ACF/PACF signatures}

\begin{center}
\begin{tabular}{lcc}
\toprule
\textbf{Model} & \textbf{ACF} & \textbf{PACF} \\
\midrule
AR$(p)$ & Exponential/sinusoidal decay & Truncates after $p$ \\
MA$(q)$ & Truncates after $q$ & Exponential/sinusoidal decay \\
ARMA$(p,q)$ & Decay after $q$ & Decay after $p$ \\
\bottomrule
\end{tabular}
\end{center}

\section{Implementation}

\begin{codeblock}[title=\textbf{Python}]
\begin{minted}{python}
import numpy as np
import matplotlib.pyplot as plt
from statsmodels.tsa.arima_process import ArmaProcess
from statsmodels.graphics.tsaplots import plot_acf, plot_pacf

# Simuler un AR(2)
np.random.seed(42)
ar_params = np.array([1, -0.75, 0.25])  # 1 - 0.75*B + 0.25*B^2
ma_params = np.array([1])
ar2 = ArmaProcess(ar_params, ma_params)
x = ar2.generate_sample(nsample=500)

fig, axes = plt.subplots(1, 3, figsize=(15, 4))
axes[0].plot(x, lw=0.7)
axes[0].set_title('AR(2): $\\phi_1=0.75, \\phi_2=-0.25$')
plot_acf(x, lags=25, ax=axes[1], title='ACF')
plot_pacf(x, lags=25, ax=axes[2], title='PACF')
plt.tight_layout()
plt.savefig('ch03_arma.pdf')
plt.show()

# Simuler un MA(2)
ma_params = np.array([1, 0.6, -0.3])
ma2 = ArmaProcess(np.array([1]), ma_params)
y = ma2.generate_sample(nsample=500)

# Simuler un ARMA(1,1)
arma11 = ArmaProcess(np.array([1, -0.8]), np.array([1, 0.4]))
z = arma11.generate_sample(nsample=500)

# Yule-Walker
from statsmodels.regression.linear_model import yule_walker
rho, sigma = yule_walker(x, order=2)
print(f"Yule-Walker: phi = {rho}")
\end{minted}
\end{codeblock}

\section{Exercises}

\begin{exercise}[$\star$]
Compute the ACF of an MA(1): $X_t = \eps_t+0.5\eps_{t-1}$. Verify numerically.
\end{exercise}

\begin{exercise}[$\star\star$]
Show that the AR(1) $X_t = \phi X_{t-1}+\eps_t$ with $|\phi|<1$ has an MA$(\infty)$ representation: $X_t = \sum_{j=0}^{\infty}\phi^j\eps_{t-j}$.
\end{exercise}

\begin{exercise}[$\star\star$]
Show that an ARMA(1,1) with $\phi=\theta$ reduces to a white noise (common root).
\end{exercise}

\begin{exercise}[$\star\star$]
Derive the Yule--Walker equations for an AR(2) and solve the system.
\end{exercise}

\begin{exercise}[$\star\star\star$ -- Project]
Simulate AR(1), AR(2), MA(1), MA(2), ARMA(1,1), and ARMA(2,1) processes. For each, plot the theoretical and empirical ACF and PACF. Build a visual identification guide.
\end{exercise}

\begin{keyformulas}
\begin{align*}
\text{AR}(p) &: \phi(B)X_t=\eps_t, \quad \text{stationary iff roots of }\phi(z)\text{ outside disc} \\
\text{MA}(q) &: X_t=\theta(B)\eps_t, \quad \text{always stationary} \\
\text{ACF AR}(1) &: \rho(h)=\phi^{|h|} \\
\text{ACF MA}(q) &: \rho(h)=0\text{ for }|h|>q
\end{align*}
\end{keyformulas}
